% Acrónimos y abreviaturas
\newacronym{api}{API}{Application Programming Interface}
\newacronym{apirest}{API REST}{Representational State Transfer Application Programming Interface}
\newacronym{sql}{SQL}{Structured Query Language}
\newacronym{html}{HTML}{HyperText Markup Language}
\newacronym{json}{JSON}{JavaScript Object Notation}

% Glosario

\newglossaryentry{backend}{
    name=Backend,
    description={Parte del sistema que procesa la lógica del servidor}
}

\newglossaryentry{reverseproxy}{
    name=Reverse Proxy (Proxy Inverso),
    description={Un tipo de servidor que recupera recursos en nombre de un cliente desde uno o más servidores back-end. Protege la identidad de los servidores internos y mejora la seguridad y el rendimiento mediante el almacenamiento en caché. \textcite{cloudflare_proxy}}
}

\newglossaryentry{framework}{
    name=Framework (Marco de Trabajo),
    description={Un conjunto estandarizado de conceptos, prácticas y criterios para enfocar un tipo de problemática particular, que sirve como referencia para enfrentar y resolver nuevos problemas de índole similar. A diferencia de una biblioteca, el framework controla el flujo de la ejecución (inversión de control). \textcite{ibm_framework}}
}

\newglossaryentry{runtime}{
    name=Runtime (Entorno de Ejecución),
    description={El sistema de software diseñado para ejecutar otros programas. Proporciona los servicios básicos y las bibliotecas necesarias para que el código escrito en un lenguaje específico (como JavaScript) pueda comunicarse con el hardware o el sistema operativo del servidor. \textcite{mdn_runtime}}
}

\newglossaryentry{API}{
    name=Runtime (API (Application Programming Interface)),
    description={Un conjunto de definiciones y protocolos que permiten que dos aplicaciones de software se comuniquen entre sí. En el desarrollo web, suele referirse a los puntos de acceso (endpoints) que un servidor expone para recibir y enviar datos en formato JSON. \textcite{aws_api}}
}

\newglossaryentry{SPA}{
    name=Runtime (SPA (Single Page Application)),
    description={Una aplicación web o sitio web que interactúa con el usuario reescribiendo dinámicamente la página actual, en lugar de cargar páginas enteras nuevas desde el servidor. Esto proporciona una experiencia de usuario más fluida, similar a una aplicación nativa. \textcite{oracle_spa}}
}


